% Fil rouge de la thèse : boucle d'optimisation en ligne
% (couleurs jvteal / jvbleu / jvgris définies dans style/presentation.cls)
\begin{tikzpicture}[
        >=latex,
        etape/.style={draw=jvteal, fill=jvteal!7, rounded corners=2pt, thick,
                      align=center, minimum width=31mm, minimum height=9.5mm,
                      font=\footnotesize\bfseries, text=jvbleu},
        coeur/.style={etape, draw=jvbleu, fill=jvbleu!8},
        procede/.style={draw=black!35, dashed, fill=black!3, rounded corners=2pt,
                      align=center, minimum width=31mm, minimum height=9.5mm,
                      font=\footnotesize, text=black!60},
        pastille/.style={circle, fill=#1, text=white, font=\tiny\bfseries,
                      inner sep=0pt, minimum size=3.8mm},
        flux/.style={->, thick, draw=jvgris},
        option/.style={draw=black!30, dashed, fill=black!4, rounded corners=2pt,
                      align=center, minimum width=31mm, minimum height=7.5mm,
                      font=\footnotesize\bfseries, text=black!45},
        fluxoption/.style={->, thick, dashed, draw=black!30},
        etiqoption/.style={etiq, text=black!40},
        etiq/.style={font=\scriptsize, text=jvgris, inner sep=1pt},
    ]
    \newcommand{\soustitre}[1]{\\[-0.4ex]{\normalfont\scriptsize\color{jvgris}#1}}
    % --- grille 3 colonnes x 2 rangées ---
    \def\dx{43mm}\def\dy{14mm}
    % rangée haute : commande
    \node[coeur]   (optim)   at (0,0)       {Optimisation\soustitre{RPEM en ligne}};
    \node[coeur]   (broche)  at (\dx,0)     {Broche asservie\soustitre{identification en boucle fermée}};
    \node[procede] (usinage) at (2*\dx,0)   {Usinage\soustitre{machine-outil}};
    % piste non traitée dans la thèse (en grisé) : avance asservie en parallèle de la broche
    \node[option]  (avance)  at (\dx,11mm)  {Moteur d'avance asservi\\[-0.4ex]{\normalfont\scriptsize perspective}};
    % rangée basse : mesure et modèle
    \node[etape]   (modele)  at (\dx,-\dy)  {Modèle de puissance\soustitre{torseurs, $K_c$, $K_f$, $K_m$}};
    \node[etape]   (instrum) at (2*\dx,-\dy){Instrumentation\soustitre{platine étalonnée par cobot}};

    % --- pastilles de section ---
    \foreach \n/\s/\c in {optim/sec:optim/jvbleu, broche/sec:ident/jvbleu,
                          modele/sec:modele/jvteal, instrum/sec:moyen_exp/jvteal}
        \node[pastille=\c] at ([xshift=0.6mm, yshift=-0.6mm]\n.north west) {\ref{\s}};

    % --- flux ---
    \draw[flux] (optim) -- node[etiq, above] {$V_c$} (broche);
    \draw[flux] (broche) -- node[etiq, above] {$N$} (usinage);
    \draw[fluxoption] (optim.north) |- node[etiqoption, pos=0.75, above] {$f_z$} (avance.west);
    \draw[fluxoption] (avance.east) -| node[etiqoption, pos=0.25, above] {$V_f$} (usinage.north);
    \draw[flux] (usinage) -- node[etiq, right] {$F,\ \mathcal{P}$} (instrum);
    \draw[flux] (instrum) -- node[etiq, above] {efforts} (modele);
    \draw[flux] (modele.west) -| node[etiq, pos=0.75, left, align=center] {énergie\\spécifique} (optim.south);
\end{tikzpicture}
